% Chapter 6 -- the regional register. Update with registers/* (manager.py, db.py, model.py, rna.py, gleif.py,
% opencoesione.py, agcm.py, lists.py, bought.py, italy_codes.py, discovery_source.py, texts.py),
% gui/register_manager.py, gui/register_settings.py, gui/window_registers.py, discovery/imports.py.
\chapter{The regional register}\label{ch:register}

Open map data finds many places but knows little about them. The \textbf{regional
register} is a list of companies kept on your PC, built from free official open data,
from the lists you import or buy, and from the companies you add. It is
\emph{one} register: official lists, imported lists and your own corrections live in the same
SQLite file (\code{<data>/registers/register.sqlite}), and the search reads it as the source
\code{regional\_register}, locally and without the network. \textbf{Nothing is downloaded until
you ask}, in \ui{Search $>$ Register manager\ldots} (also a button in tab~2).

\section{The four official sources}\label{sec:regsources}
All four are on by default; each keeps its licence and attribution (shown in the manager, in the
\ui{Licence} column). Only companies are stored: natural persons (16-character tax
codes, sole traders, professionals) are not.

{\small
\begin{tabularx}{\textwidth}{@{}l l Y@{}}
\toprule
\hdr \thead{Source} & \thead{Licence} & \thead{What it gives, size, refresh} \\
\midrule
RNA, national register of state aid (Incentivi Italia) & CC BY 4.0 & Companies that received state aid, de minimis aid, guarantees or tax credits: name, \emph{region only}, size class, a sector \emph{hint}, number and date of aids, hiring incentives as a signal. Monthly ZIP files with XML up to about 600\,MB; the last 24 months by default (Settings: 0--120, 0 = whole archive, 1--2\,GB). Refreshed monthly. \\
GLEIF golden copy (LEI register) & CC0 1.0 & Companies with an LEI: legal name, registered address, postcode, municipality, province, legal form, founding date, status; the tax code when it is an 11-digit number. About 490\,MB zipped (5\,GB of CSV worldwide), read from the ZIP without extracting it. Treated as monthly. \\
OpenCoesione, recipients of cohesion funds & CC BY 4.0 & Municipality, address, the \emph{full 6-digit ATECO code} and its description, employee and size class; public bodies excluded. One CSV of about 2.3\,GB in a ZIP. About twice a year. \\
AGCM, legality rating & CC BY 4.0 & Companies with a legality rating: stars, validity, municipality and province, town-level coordinates. Weekly list; the site's 10-second crawl delay is respected. \emph{Statistical value only}: later revocations must be checked on newer lists. \\
\bottomrule
\end{tabularx}
}

\vspace{3pt}
Downloads go through the polite gateway (2\,s between requests), resume after an
interruption (a \code{.part} file and an HTTP range), are checked (size, ZIP structure,
SHA-256 where published), and a damaged file is discarded while the previous good copy is kept.
\ui{Check for updates} asks each source whether it changed and how much would be fetched
(no download); \ui{Download / update} asks first, and its result says which RNA months were loaded (``Months loaded: 12/2024--11/2026 (24)'') and which were not yet published and skipped; \ui{Stop} keeps the partial files. A refresh
replaces, per source and region, that source's earlier values; a company no source lists any more
disappears, unless you added or corrected it. Settings, \ui{Registers and paid key}, chooses the
\emph{regions to keep} (default Emilia-Romagna), the RNA months and the sources, and can let the
program \emph{check} (never download) about once a month at start-up.

A note on coverage (how many companies the free registers list for the province chosen in tab~2 and roughly what share of the active ones that is) is shown in tab~2, under the register's controls, and under the table of the Register manager.

\section{Merge, provenance, your corrections}
\textbf{Merge by tax code.} The key is the 11-digit tax code (an ``IT'' prefix is removed,
9--10 digits are zero-padded). Rows with the same code are one company; the same name with
another code stays separate, and \ui{Possible duplicates} lists similar names without ever
merging them. A row with no code (from a list or added by hand) gets a stable internal key.

\textbf{Per-field provenance.} The register keeps \emph{every} source's value with its
source and date and shows one by a fixed priority per field: you first, then the list you
bought, then for example GLEIF for name and address, OpenCoesione for the exact ATECO code,
AGCM for coordinates; your imported lists only when no official source has the field. The
\ui{Sources} column names where a row comes from; a \ui{Signals} column shows badges (``Legality
rating 3 stars'', ``Received public aid'', ``Hired with incentives''), and company cards get the
same facts with the source named. An ATECO known only from RNA stays a \emph{hint}
(``indicative: sector declared for an aid''), never a code.

\textbf{Your corrections are locked.} Double-click a cell to change it (name, VAT, place,
province, ATECO, website, legal form, employees, status, e-mail, phone): it gets a
padlock and a refresh \emph{never} overwrites it; \ui{Unlock field} brings the sources' value
back. You can \ui{Add company}, \ui{Hide} one (it stays hidden after every refresh),
\ui{Delete} (a company you added is removed, one a source lists is hidden) and \ui{Undo} up to
200 steps. \ui{Export CSV} writes the whole register.

\section{Province, approximated and flagged}
When no source gives a province and you have not set one, the program reads digits 8--10 of
the tax code, the \emph{tax office of first registration}, and maps them to a province. Where
a code is shared by provinces created later, candidates are filtered by the company's region;
one left is used and marked \emph{presumed} (tooltip ``from the tax office code''), several
leave the province empty with the candidates kept. A search in one province includes presumed
ones. Region and province otherwise come from the postcode (GLEIF, bought lists) or the ISTAT
municipality code (OpenCoesione).

\section{Imported and bought lists}\label{sec:lists}
\ui{Import list (CSV/XLSX)} adds any company list, with a preview to map the columns;
\ui{Import bought Registro Imprese list} reads the ``Elenco'' you bought there, whose values rank
just after your own corrections. Each list is a source of its own: \ui{Refresh list from file}
replaces it with a newer file and keeps your corrections, \ui{Delete list} removes it (asks first)
while companies another source lists stay. Lists kept by older versions are moved into the
register once; the old copy stays as a backup.

\section{How the search uses it, and how much it covers}
For a province or region search the register supplies the companies of that area. The ATECO
categories you chose select by exact code and by RNA hint, companies with no code are kept when
they come from a list or from you, hidden and inactive ones are left out, and at most 3\,000 are
read (the status line says ``choose ATECO categories to narrow''). Position comes from AGCM or
from the town in the local map data; a radius search drops companies it cannot place. The
register and your lists are the only trusted source of a company's exact ATECO code and legal
name. The pre-screen (chapter~\ref{ch:interest}) and the second job-search list
(chapter~\ref{ch:jobs}) read it too.

\begin{warnbox}
\textbf{The free registers are not complete.} In the project's own check on the machinery
companies of one province they listed roughly a quarter to a third of the companies the
Chambers of Commerce count; website, e-mail, phone and PEC are not in them at all. The
complete list is the \emph{paid} Registro Imprese ``Elenco'' (5\,EUR plus 0.02--0.12\,EUR per
company), which you buy yourself and import. Read a missing company as ``not in the lists'',
not as ``does not exist''.
\end{warnbox}
